% GENERATED FILE. Edit canonical lessons, then run npm run generate:books.
% canonical-source-hash: fnv1a64:77a3e196e7212300
% canonical-lessons: IT-C418-biglietteria, IT-C418-abbonamento, IT-C418-timbrare, IT-C418-controllore, IT-C418-supplemento

\chapter{Il biglietto}
\label{ch:it-il-biglietto}
\hlchaptermodality{418}

\begin{chapteropening}
\textbf{By the end of this chapter:} \emph{I can buy and use tickets: the ticket office, a season ticket, stamping it, the inspector and the supplement.}

Complete the last lesson of chapter 418: I can buy and use tickets: the ticket office, a season ticket, stamping it, the inspector and the supplement.
\end{chapteropening}

\section[la biglietteria]{la biglietteria — the ticket office}

\label{lesson:IT-C418-biglietteria}



\begin{quote}
Before the new one: say the Italian for the long-distance bus, then the Italian for to pass through.
\end{quote}

\subsection*{You'll want to know: la biglietteria}
\textbf{la biglietteria} — \textquotedblleft{}the ticket office\textquotedblright{}.

\begin{quote}
La biglietteria apre alle sei.
\end{quote}

\begin{quote}
\emph{The ticket office opens at six.}
\end{quote}

\emph{La biglietteria automatica} is the ticket machine.

\begin{cousinweb}
\emph{biglietteria} is built on \emph{biglietto} 'ticket', which was borrowed from French \emph{billet} 'note' (EtymDB).
\end{cousinweb}

\begin{grammarlens}[title={What you've built}]
One more word for telling, explaining and describing at B1.
\end{grammarlens}

\subsection*{Your turn}
\begin{itemize}
\raggedright
  \item \emph{Say it:} \emph{la biglietteria}
  \item \emph{Say it:} \emph{la biglietteria}, once more
  \item \emph{Say it:} a sentence of your own with \emph{la biglietteria}
  \item \emph{Recall:} say the Italian for the long-distance bus, then the Italian for to pass through, then say \emph{la biglietteria} again
\end{itemize}

\begin{tcolorbox}[breakable,colback=teal!4,colframe=teal!35!black,title={Before you move on}]
What does la biglietteria mean? (\textquotedblleft{}The ticket office\textquotedblright{}.) Say it once more.
\end{tcolorbox}
\section[l'abbonamento]{l'abbonamento — the season ticket, the pass; the subscription}

\label{lesson:IT-C418-abbonamento}



\begin{quote}
Before the new one: say the Italian for to pass through, then the Italian for the ticket office.
\end{quote}

\subsection*{You'll want to know: l'abbonamento}
\textbf{l'abbonamento} — \textquotedblleft{}the season ticket, the pass; the subscription\textquotedblright{}.

\begin{quote}
Ho l'abbonamento mensile per il tram.
\end{quote}

\begin{quote}
\emph{I have a monthly pass for the tram.}
\end{quote}

For a newspaper or a gym too.

\begin{cousinweb}
\emph{abbonamento} is built on \emph{abbonare} 'to subscribe', which was borrowed from French \emph{abonner} (EtymDB).
\end{cousinweb}

\begin{grammarlens}[title={What you've built}]
One more word for telling, explaining and describing at B1.
\end{grammarlens}

\subsection*{Your turn}
\begin{itemize}
\raggedright
  \item \emph{Say it:} \emph{l'abbonamento}
  \item \emph{Say it:} \emph{l'abbonamento}, once more
  \item \emph{Say it:} a sentence of your own with \emph{l'abbonamento}
  \item \emph{Recall:} say the Italian for to pass through, then the Italian for the ticket office, then say \emph{l'abbonamento} again
\end{itemize}

\begin{tcolorbox}[breakable,colback=teal!4,colframe=teal!35!black,title={Before you move on}]
What does l'abbonamento mean? (\textquotedblleft{}The season ticket, the pass; the subscription\textquotedblright{}.) Say it once more.
\end{tcolorbox}
\section[timbrare]{timbrare — to stamp, to validate (a ticket)}

\label{lesson:IT-C418-timbrare}



\begin{quote}
Before the new one: say the Italian for the ticket office, then the Italian for the season ticket.
\end{quote}

\subsection*{You'll want to know: timbrare}
\textbf{timbrare} — \textquotedblleft{}to stamp, to validate (a ticket)\textquotedblright{}.

\begin{quote}
Ricordati di timbrare il biglietto prima di salire.
\end{quote}

\begin{quote}
\emph{Remember to stamp your ticket before you get on.}
\end{quote}

At work, \emph{timbrare il cartellino} is to clock in.

\begin{cousinweb}
\emph{timbrare} is built on \emph{timbro} 'stamp', which was borrowed from French \emph{timbre}, a word that goes back through Latin \emph{tympanum} to Greek \emph{týmpanon} 'drum'.
\end{cousinweb}

\begin{grammarlens}[title={What you've built}]
One more word for telling, explaining and describing at B1.
\end{grammarlens}

\subsection*{Your turn}
\begin{itemize}
\raggedright
  \item \emph{Say it:} \emph{timbrare}
  \item \emph{Say it:} \emph{timbrare}, once more
  \item \emph{Say it:} a sentence of your own with \emph{timbrare}
  \item \emph{Recall:} say the Italian for the ticket office, then the Italian for the season ticket, then say \emph{timbrare} again
\end{itemize}

\begin{tcolorbox}[breakable,colback=teal!4,colframe=teal!35!black,title={Before you move on}]
What does timbrare mean? (\textquotedblleft{}To stamp, to validate (a ticket)\textquotedblright{}.) Say it once more.
\end{tcolorbox}
\section[il controllore]{il controllore — the ticket inspector}

\label{lesson:IT-C418-controllore}



\begin{quote}
Before the new one: say the Italian for the season ticket, then the Italian for to stamp.
\end{quote}

\subsection*{You'll want to know: il controllore}
\textbf{il controllore} — \textquotedblleft{}the ticket inspector\textquotedblright{}.

\begin{quote}
Il controllore passa tra i sedili.
\end{quote}

\begin{quote}
\emph{The inspector comes along the seats.}
\end{quote}

\emph{Il controllore di volo} is an air traffic controller.

\begin{cousinweb}
\emph{controllore} is built on \emph{controllo} 'check', which was borrowed from French \emph{contrôle}, from \emph{contre-rôle}: a second list kept to check the first.
\end{cousinweb}

\begin{grammarlens}[title={What you've built}]
One more word for telling, explaining and describing at B1.
\end{grammarlens}

\subsection*{Your turn}
\begin{itemize}
\raggedright
  \item \emph{Say it:} \emph{il controllore}
  \item \emph{Say it:} \emph{il controllore}, once more
  \item \emph{Say it:} a sentence of your own with \emph{il controllore}
  \item \emph{Recall:} say the Italian for the season ticket, then the Italian for to stamp, then say \emph{il controllore} again
\end{itemize}

\begin{tcolorbox}[breakable,colback=teal!4,colframe=teal!35!black,title={Before you move on}]
What does il controllore mean? (\textquotedblleft{}The ticket inspector\textquotedblright{}.) Say it once more.
\end{tcolorbox}
\section[il supplemento]{il supplemento — the supplement, the extra charge}

\label{lesson:IT-C418-supplemento}



\begin{quote}
Before the new one: say the Italian for to stamp, then the Italian for the ticket inspector.
\end{quote}

\subsection*{You'll want to know: il supplemento}
\textbf{il supplemento} — \textquotedblleft{}the supplement, the extra charge\textquotedblright{}.

\begin{quote}
Per il treno veloce c'è un supplemento.
\end{quote}

\begin{quote}
\emph{For the fast train there is a supplement.}
\end{quote}

Also a newspaper supplement.

\begin{cousinweb}
\emph{supplemento} is a learned borrowing from Latin \emph{supplēmentum} 'a filling up' (EtymDB), from \emph{supplēre} 'to fill up', built on \emph{plēnus} 'full'.
\end{cousinweb}

\begin{grammarlens}[title={What you've built}]
One more word for telling, explaining and describing at B1.
\end{grammarlens}

\subsection*{Your turn}
\begin{itemize}
\raggedright
  \item \emph{Say it:} \emph{il supplemento}
  \item \emph{Say it:} \emph{il supplemento}, once more
  \item \emph{Say it:} a sentence of your own with \emph{il supplemento}
  \item \emph{Recall:} say the Italian for to stamp, then the Italian for the ticket inspector, then say \emph{il supplemento} again
\end{itemize}

\begin{tcolorbox}[breakable,colback=teal!4,colframe=teal!35!black,title={Before you move on}]
What does il supplemento mean? (\textquotedblleft{}The supplement, the extra charge\textquotedblright{}.) Say it once more.
\end{tcolorbox}
